\begin{scanpage}{0033}{16}
Here the absolute value of the lower limit of the Fresnel integral always equals infinity, and its sign is determined by the sign of $\cos\psi/2$. Therefore, when passing through the boundary of the plane waves $(\psi=\phi\pm\phi_0=\pi)$ the lower limit changes sign and the Fresnel integral undergoes a finite discontinuity, ensuring at this boundary the continuity of the function $u(r,\psi)$ and consequently of the diffraction field. Actually, by means of the well-known equation
\[\int_0^\infty e^{iq^2}\,dq=\frac{\sqrt\pi}{2}e^{i\frac\pi4}\tag{2.14}\]
it is not difficult to show that
\[v(r,\pi+0)=\frac{e^{ikr}}2,\quad v(r,\pi-0)=-\frac{e^{ikr}}2\tag{2.15}\]
and consequently
\[u(r,\pi\pm0)=\frac12e^{ikr}.\tag{2.16}\]
In view of the asymptotic relationships
\[\int_\infty^p e^{iq^2}\,dq=\frac{e^{ip^2}}{2ip},\quad\int_{-\infty}^{-p}e^{iq^2}\,dq=-\frac{e^{ip^2}}{2ip}\quad(\text{with }p\gg1)\tag{2.17}\]
Pauli's Equation (2.13) is transformed with $\sqrt{2kr}\left|\cos\frac\psi2\right|\gg1$ to the Expression (2.04). As was already indicated above, it determines the cylindrical waves diverging from the wedge edge.

By means of Equation (2.13), one may also calculate the field in the vicinity of the direction $\phi=2\alpha-\pi-\phi_0$ --- that is, near the boundary of the plane wave reflected from the face $\phi=\alpha$; for this purpose, it is sufficient to replace $\phi$ by $\alpha-\phi$ and $\phi_0$ by $\alpha-\phi_0$.

It is also interesting to note that in the case of a half-plane $(n=2)$ Equation (2.13) gives the expression
\end{scanpage}
